\documentclass[letterpaper,10.8pt]{article}

\usepackage{latexsym}
\usepackage[empty]{fullpage}
\usepackage{titlesec}
\usepackage{marvosym}
\usepackage[usenames,dvipsnames]{color}
\usepackage{verbatim} 
\usepackage{enumitem}
\usepackage[pdftex]{hyperref}
\usepackage{fancyhdr}

% \usepackage[hidelinks]{hyperref}
\usepackage{fontawesome5}

\pagestyle{fancy}
\fancyhf{} % clear all header and footer fields
\fancyfoot{}
\renewcommand{\headrulewidth}{0pt}
\renewcommand{\footrulewidth}{0pt}

% Adjust margins
\addtolength{\oddsidemargin}{-0.375in}
\addtolength{\evensidemargin}{-0.375in}
\addtolength{\textwidth}{1in}
\addtolength{\topmargin}{-.8in}
\addtolength{\textheight}{1.5in}

\urlstyle{tt}

\raggedbottom
\raggedright
\setlength{\tabcolsep}{0in}

% Sections formatting
\titleformat{\section}{
  \vspace{-3pt}\scshape\raggedright\large
}{}{0em}{}[\color{black}\titlerule \vspace{-5pt}]

%-------------------------
% Custom commands
\newcommand{\resumeItem}[2]{
  \item\small{
    \textbf{#1}{ #2 }
  }
}

% Publications: justified text, [Paper] link kept in one piece
\newcommand{\resumeItemPublication}[2]{
  \item\footnotesize{\rightskip=0pt\relax\emergencystretch=3em\leavevmode
    #1\unskip~\mbox{#2}\par
  }
}

\newcommand{\venue}[1]{\textls[15]{\textsc{\textcolor{gray!70}{#1}}}}
\newcommand{\journal}[1]{\textls[15]{\textsc{\textcolor{gray!70}{#1}}}}

\newcommand{\resumeItemWithoutTitle}[1]{
  \item\small{
    {}
  }
}

\newcommand{\resumeSubheading}[4]{
  \vspace{-1pt}\item[]
    \textbf{#1}\hfill #2\\[1pt]
    \textit{\small#3}\hfill\textit{\small #4}\par\vspace{2pt}%
}

\newcommand{\resumeSubheadingTwo}[2]{
  \vspace{0pt}\item[]
    \textbf{#1}\hfill #2\par\nopagebreak\vspace{1pt}%
}

\newcommand{\resumeSubItem}[2]{\resumeItem{#1}{#2}}

% Uniform filled bullets at every level
\renewcommand{\labelitemi}{$\bullet$}
\renewcommand{\labelitemii}{$\bullet$}

\newcommand{\resumeSubHeadingListStart}{\begin{itemize}[leftmargin=0pt, labelwidth=0pt, labelsep=0pt, labelindent=0pt, itemindent=0pt]}
\newcommand{\resumeSubHeadingListEnd}{\end{itemize}}

\newcommand{\resumeSubHeadingNumberStart}{%
  \begin{enumerate}[leftmargin=2.8em, labelindent=0.4em, labelwidth=1.7em, labelsep=0.7em, align=left, label={[\arabic*]}]%
}
\newcommand{\resumeSubHeadingNumberEnd}{\end{enumerate}}

\newcommand{\resumeItemListStart}{\begin{itemize}}
\newcommand{\resumeItemListEnd}{\end{itemize}\vspace{-5pt}}

%-------------------------------------------
%%%%%%  CV STARTS HERE  %%%%%%%%%%%%%%%%%%%%%%%%%%%%

\begin{document}
\small
\setlist[itemize]{leftmargin=*,labelindent=0.6em,labelsep=0.5em,topsep=3pt,parsep=0pt,partopsep=0pt}

%----------HEADING-----------------
\begin{center}
  {\LARGE \textbf{VIBHAS REDDY ARAVEETI}}\\[4pt]
  \faGlobe\;\href{https://portfolio-vibhasf17-4811.vercel.app/}{Portfolio} \quad
  \faEnvelope\;\href{mailto:vibhasf17@gmail.com}{vibhasf17@gmail.com} \quad
  \faPhone\;+91 63013 26306 \quad
  \faLinkedin\;\href{https://www.linkedin.com/in/araveeti-vibhas-reddy-329b19291/}{LinkedIn} \quad
  \faGraduationCap\;\href{https://scholar.google.com/citations?hl=en&view_op=list_works&gmla=AN_4Zu8eJn7hGADizn1nlrwoGJPhVHfqSGZ8C3mjxd3F5X6qqm9l9fAdWKvJLqJBGBSIPndmTOc00tn0eIFiyw&user=d_uIop8AAAAJ}{Google Scholar}
\end{center}

%-----------EDUCATION-----------------
\section{Education}
  \resumeSubHeadingListStart
    \resumeSubheading
      {Amrita Vishwa Vidyapeetham}{Bengaluru, India}
      {Bachelor of Technology, Computer Science and Engineering, \textbf{GPA: 8.65/10}}{2023 -- 2027}
    \item[] \small\textbf{Relevant Coursework:} Data Structures and Algorithms, Object-Oriented Programming, Operating Systems, Computer Organization, Database Management Systems, Computer Networks, Machine learning, Deep learning, natural language processing, Distributed systems, Probability, Calculus, Linear Algebra, Discrete mathematics
  \resumeSubHeadingListEnd

%-----------PUBLICATIONS-----------------
\section{Publications}
  \resumeSubHeadingNumberStart
  \setlength{\itemsep}{0pt}\setlength{\parsep}{0pt}
  \resumeItemPublication{\llap{``}Heart Disease Risk Prediction Using Deep Learning and Probabilistic Calibration Techniques,'' \textit{Engineered Science}, Q1 journal.}
  {\textbf{\href{https://www.espublisher.com/journals/articledetails/2461}{[Paper]}}}

  \resumeItemPublication{\llap{``}AI-Powered Waste Detection and Sorting Using Object Detection,'' in \textit{Proc. IEEE Int. Conf. Self-Sustainable Artificial Intelligence Systems (ICSSAS)}, 2026.}
  {\textbf{\href{https://ieeexplore.ieee.org/document/11559513}{[Paper]}}}

  \resumeItemPublication{\llap{``}Detecting Health Misinformation: AI Solutions for Medical and Policy Content,'' in \textit{Proc. IEEE Int. Conf. Knowledge Engineering and Communication Systems (ICKECS)}, 2026. \textbf{\href{https://github.com/VIBHAS777/research-papers/blob/main/13-Detecting-Health-Misinformation-AI-Solutions/Best_Paper_Award_ICKECS_2026.pdf}{Best Paper Award.}}}
  {\textbf{\href{https://ieeexplore.ieee.org/document/11528010}{[Paper]}}}

  \resumeItemPublication{\llap{``}Efficient Malayalam Handwritten Character Recognition Using ConvNeXt-PCA-RandomForest Architecture,'' in \textit{Proc. IEEE Int. Conf. Convergence in Computing Technology (I3CTCON)}, 2026.}
  {\textbf{\href{https://ieeexplore.ieee.org/document/11507971}{[Paper]}}}

  \resumeItemPublication{\llap{``}Enhancing CPU Benchmark Scores Using RNN-Based Deadlock Detection from System Logs,'' presented at the 16th ICCCNT, 2025.}
  {\textbf{\href{https://github.com/VIBHAS777/research-papers/tree/main/05-RNN-Based-Deadlock-Detection}{[Paper]}}}

  \resumeItemPublication{\llap{``}Optimization of Space Shuttle External Tank Design Return on Investment (ROI),'' presented at ICAANA, 2025.}
  {\textbf{\href{https://github.com/VIBHAS777/research-papers/blob/main/06-Space-Shuttle-External-Tank-ROI/Optimization_of_Space_Shuttle_External_Tank_Design_ROI.pdf}{[Paper]}}}

  \resumeItemPublication{\llap{``}Automated Medical Image Analysis for Brain Tumor and Skin Disease Diagnosis Using Transfer Learning and Grad-CAM,'' presented at the 3rd ICITEICS, 2026.}
  {\textbf{\href{https://github.com/VIBHAS777/research-papers/tree/main/07-Automated-Medical-Image-Analysis-Grad-CAM}{[Paper]}}}

  \resumeItemPublication{\llap{``}AI-Driven Load Balancer with Blockchain-Based Decision Making in Banking System,'' presented at the 10th ICECA, 2026.}
  {\textbf{\href{https://github.com/VIBHAS777/research-papers/tree/main/08-AI-Driven-Load-Balancer-Banking}{[Paper]}}}

  \resumeItemPublication{\llap{``}Stability Threshold and Performance Degradation Analysis of Priority-Aware SDN in Healthcare IoT Systems,'' presented at ICUIS, 2026.}
  {\textbf{\href{https://github.com/VIBHAS777/research-papers/tree/main/09-Priority-Aware-SDN-Healthcare-IoT}{[Paper]}}}

  \resumeItemPublication{\llap{``}Intelligent Congestion Prediction and Control in Data Center Networks Using SDN and Deep Learning,'' accepted at ICON INDIA, 2026.}
  {\textbf{\href{https://github.com/VIBHAS777/research-papers/tree/main/10-SDN-Deep-Learning-Congestion-Control}{[Paper]}}}

  \resumeItemPublication{\llap{``}A Unified Machine Learning Pipeline for Early Prediction of Heart Disease, Diabetes Mellitus, and Alzheimer's Disease,'' accepted at ICRTAC, 2026.}
  {\textbf{\href{https://github.com/VIBHAS777/research-papers/tree/main/11-Unified-Chronic-Disease-Prediction-Pipeline}{[Paper]}}}

  \resumeItemPublication{\llap{``}A Low-Cost Real-Time Multimodal Seizure Detection System Using MAX30102 and Accelerometry,'' under review at ICACS, 2026.}
  {\textbf{\href{https://github.com/VIBHAS777/research-papers/blob/main/12-Multimodal-Seizure-Detection-MAX30102/Low_Cost_Real_Time_Multimodal_Seizure_Detection_System_MAX30102.pdf}{[Paper]}}}
  \resumeSubHeadingNumberEnd

%-----------EXPERIENCE-----------------
\section{Work Experience}
  \resumeSubHeadingListStart
    \resumeSubheadingTwo
    {Amrita University}{Bengaluru, India}

    \begin{itemize}[leftmargin=0pt, labelwidth=0pt, labelsep=0pt, labelindent=0pt, itemindent=0pt, itemsep=0pt]
            \item[] \small\textbf{Research Intern} \hfill \textit{Mar 2025 -- Jun 2026}
        \begin{itemize}[leftmargin=*, itemsep=0.5pt, label={$\bullet$}]
          \item Investigated heart attack risk prediction to improve the reliability of clinical risk assessment through deep learning and probabilistic calibration, under the supervision of Dr. B. Venkatesh.
          \item Built and trained a bagged Multilayer Perceptron (MLP) ensemble framework in Python using scikit-learn and PyTorch, integrating Bayesian bootstrap aggregation with multi-technique probability calibration (Platt scaling, isotonic regression, and histogram binning).
          \item Engineered domain-specific clinical preprocessing pipelines and interaction features, including Cholesterol $\times$ BMI, pulse pressure to BP ratio, and Age $\times$ Stress level, using Pandas and NumPy.
          \item Achieved 96.42\% model accuracy with a Brier score of 0.0489, a recall of 0.97, and a false negative rate reduced to 3\% through rigorous data preparation and calibration, ensuring predicted risk scores reflected real-world medical frequencies.
          \item Co-authored the research, published in the Q1-indexed journal Engineered Science (DOI: 10.30919/es2461).
        \end{itemize}
    \end{itemize}
  \resumeSubHeadingListEnd

%-----------PROJECTS-----------------
\section{Academic Projects}
  \resumeSubHeadingListStart
    \resumeSubheadingTwo
      {Automated Medical Image Analysis for Brain Tumor and Skin Disease Diagnosis\ \ \href{https://github.com/VIBHAS777/medical-image-analysis-gradcam}{GitHub}}{\textit{Aug 2024 -- Dec 2024}}
    \begin{itemize}[leftmargin=*, itemsep=0.5pt, label={$\bullet$}]
      \item Developed a unified deep learning pipeline for brain tumor and skin disease diagnosis, applying transfer learning with VGG16, ResNet-50, and EfficientNet models to address variability across medical imaging domains.
      \item Compared CNN and transfer-learning architectures across MRI and dermoscopic datasets to evaluate classification performance across different medical imaging tasks.
      \item Improved diagnostic performance and interpretability, achieving 94\% brain tumor accuracy and 99\% skin disease accuracy with 0.993 AUC, while using Grad-CAM to identify image regions driving model predictions.
    \end{itemize}

    \resumeSubheadingTwo
      {Unified Disease Prediction Pipeline\ \ \href{https://github.com/VIBHAS777/unified-chronic-disease-prediction}{GitHub}}{\textit{Aug 2025 -- Nov 2025}}
    \begin{itemize}[leftmargin=*, itemsep=0.5pt, label={$\bullet$}]
      \item Built an end-to-end machine learning pipeline for early prediction of heart disease, diabetes, and Alzheimer's disease, using disease-specific preprocessing, SMOTE, and feature engineering across multiple classification models.
      \item Evaluated multiple classifiers across three disease-specific datasets using stratified data splitting and standardized preprocessing to compare predictive performance across clinical domains.
      \item Achieved 89.3\% accuracy for heart disease, 76.6\% for diabetes, and 82.47\% accuracy with a 0.8984 AUC-ROC for Alzheimer's, using the best-performing model for each disease.
    \end{itemize}

    \newpage
    \resumeSubheadingTwo
      {AI-Powered Health Misinformation Detection\ \ \href{https://github.com/VIBHAS777/health-misinformation-ai-detection}{GitHub}}{\textit{Oct 2025 -- Feb 2026}}
    \begin{itemize}[leftmargin=*, itemsep=0.5pt, label={$\bullet$}]
      \item Addressed class imbalance in biomedical misinformation detection by applying T5-based paraphrasing to the underrepresented Refute class and introducing an enhanced three-column claim-evidence-question representation.
      \item Compared biomedical transformer models, including BioBERT, RoBERTa, SciBERT, and ClinicalBERT, against Meta-Llama-3-8B-Instruct for three-class claim verification: Support, Refute, and Neutral.
      \item Achieved 83.91\% accuracy and 0.7831 Macro-F1 using zero-shot Meta-Llama-3-8B-Instruct, while BioBERT with augmentation improved Refute-class F1 from 0.6486 to 0.7264.
    \end{itemize}

    \resumeSubheadingTwo
      {A Low-Cost Real-Time Multimodal Seizure Detection System\ \ \href{https://github.com/VIBHAS777/multimodal-seizure-detection-max30102}{GitHub}}{\textit{Feb 2026 -- May 2026}}
    \begin{itemize}[leftmargin=*, itemsep=0.5pt, label={$\bullet$}]
      \item Built a low-cost real-time seizure monitoring prototype combining MAX30102 heart-rate sensing, MPU6050 accelerometry, and an Arduino Uno for multimodal physiological monitoring.
      \item Designed a signal fusion algorithm using physiological thresholds to cut false alarms from isolated heart-rate or motion anomalies, enabling real-time alerts with event timestamping.
      \item Integrated concurrent heart-rate and motion threshold checks to trigger audiovisual alerts when both signals indicated abnormal activity, supporting continuous monitoring outside clinical settings.
    \end{itemize}

    \resumeSubheadingTwo
      {RNN-Based Deadlock Detection for CPU Performance Optimization\ \ \href{https://github.com/VIBHAS777/rnn-deadlock-detection}{GitHub}}{\textit{Feb 2025 -- May 2025}}
    \begin{itemize}[leftmargin=*, itemsep=0.5pt, label={$\bullet$}]
      \item Designed a Python/PyTorch system that generated simulated OS logs and used a 2-layer RNN to detect impending deadlocks from sequential process and resource events.
      \item Processed simulated operating-system logs into sequential, one-hot-encoded event representations to capture temporal dependencies in process and resource allocation patterns.
      \item Achieved 90.25\% test accuracy in deadlock prediction and raised the simulated CPU benchmark score from 70 to 98.02, representing a 28.02\% improvement through proactive detection.
    \end{itemize}

    \resumeSubheadingTwo
      {Space Shuttle External Tank ROI Optimization\ \ \href{https://github.com/VIBHAS777/space-shuttle-tank-roi-optimization}{GitHub}}{\textit{Jan 2025 -- May 2025}}
    \begin{itemize}[leftmargin=*, itemsep=0.5pt, label={$\bullet$}]
      \item Developed a data-driven framework combining SLSQP constrained optimization and ML regression to optimize tank dimensions, material properties, cost, and structural constraints for ROI maximization.
      \item Applied SLSQP-based constrained optimization to determine cost-effective tank designs while maintaining structural integrity and stability under operational constraints.
      \item Achieved a strong regression performance with an R$^2$ of 0.9086, MAE of 0.0129, and MSE of 0.000451, while identifying an optimized tank configuration yielding a 44.35\% ROI under volume, stress, and vibration constraints.
    \end{itemize}

    \resumeSubheadingTwo
      {Stability Threshold and Performance Degradation Analysis of Priority-Aware SDN\ \ \href{https://github.com/VIBHAS777/priority-aware-sdn-healthcare-iot}{GitHub}}{\textit{Jan 2026 -- Jun 2026}}
    \begin{itemize}[leftmargin=*, itemsep=0.5pt, label={$\bullet$}]
      \item Engineered a priority-aware SDN framework combining Dijkstra-based dynamic routing with tabular Q-learning to adapt queue weights across four healthcare IoT traffic classes under varying network loads.
      \item Assessed adaptive traffic prioritization under varying network loads to determine its effectiveness in maintaining reliable and low-latency communication for critical healthcare IoT traffic.
      \item Reduced overall P95 latency by 40\% and jitter by 84\% at peak load, while maintaining 100\% on-time delivery for Critical traffic and lowering the Weighted Performance Degradation Index from 0.302 to 0.276.
    \end{itemize}
  \resumeSubHeadingListEnd

%-----------LEADERSHIP-----------------
\section{Leadership Experience}
  \resumeSubHeadingListStart
    \resumeSubheadingTwo
      {Amrita Vishwa Vidyapeetham}{Bengaluru, India}

    \begin{itemize}[leftmargin=0pt, labelwidth=0pt, labelsep=0pt, labelindent=0pt, itemindent=0pt, itemsep=0pt]
            \item[] \small\textbf{Co-Founder \& Lead Developer -- ML Minds} \hfill \textit{Sep 2024 -- Present}
        \begin{itemize}[leftmargin=*, itemsep=0.5pt, label={$\bullet$}]
          \item Co-founded and led a student-driven machine learning community, organizing technical workshops and initiatives for members.
          \item Mentored peer members on ML concepts and projects, fostering a collaborative learning environment within the group.
        \end{itemize}

      \vspace{2pt}
      \item[] \small\textbf{Captain -- Powerlifting Team} \hfill \textit{Jul 2024 -- Feb 2025}
        \begin{itemize}[leftmargin=*, itemsep=0.5pt, label={$\bullet$}]
          \item Guided team members to compete in competitive powerlifting events.
          \item Led the university powerlifting team, coordinating training schedules across the squad.
        \end{itemize}

      \vspace{2pt}
      \item[] \small\textbf{Cultural Coordinator -- Kalanjali 2025 (College Fest)} \hfill \textit{Apr 2025}
        \begin{itemize}[leftmargin=*, itemsep=0.5pt, label={$\bullet$}]
          \item Managed event logistics for Kalanjali 2025 and directed student teams to execute the event successfully.
        \end{itemize}
    \end{itemize}
  \resumeSubHeadingListEnd

%-----------Awards-----------------
\section{Awards and Achievements}
\begin{itemize}[leftmargin=*, itemsep=1pt, label={$\bullet$}]
\item {\small Best Paper Award, 2026 International Conference on Knowledge Engineering and Communication Systems (ICKECS 2026), for ``Detecting Health Misinformation: AI Solutions for Medical and Policy Content.''}
\item {\small Recognized among the top 3\% of student researchers across the university for research contributions.}
\item {\small Secured 1st place in the Arista Networks Boot Camp 2026 with a score of 86\%.}
\end{itemize}

%-----------Skills-----------------
\section{Skills and Interests}
	\begin{itemize}[leftmargin=*, itemsep=1pt, label={$\bullet$}]
	\resumeSubItem{Technical Skills:}{Python, Java, JavaScript, C, C++, TypeScript, Go, Rust, Ruby, SQL, NumPy, Pandas, SciPy, Scikit-learn, PyTorch, TensorFlow, Keras, LightGBM, XGBoost, CatBoost, Random Forest, SMOTE, CNNs, Transfer Learning, Reinforcement Learning, OpenCV, Grad-CAM, Explainable AI (XAI), SHAP, LIME, NiBabel, Matplotlib, Seaborn, MLflow, FHIR, HL7}
	\resumeSubItem{Developer Tools:}{VS Code, Git, GitHub, Jupyter Notebook, Docker, Kubernetes, Postman, Linux/Bash}
	\resumeSubItem{Other Interests:}{Open-Source Development, Hackathons, Competitive Programming, Technical Mentoring, Powerlifting, Swimming, Public Speaking, Aviation and Space Exploration, Social Entrepreneurship, Photography}
\end{itemize}

\end{document}